\section{What symmetry does and does not imply for fairness}\label{sec:consequences}
The common-payoff conclusion of Theorem~\ref{thm:indexing} is a strong
numerical restriction. It should not be substituted for a participant's
indifference between player assignments. The following elementary comparisons
state the extra decision data explicitly; the broader fairness programme is
not required for the structural classification.

\Needspace{11\baselineskip}
\begin{proposition}[Assignment calibration]\label{prop:assignment-calibration}
Fix conditional profile laws $L_i$ for a participant assigned to role $i$.
Within-role expected-utility preferences do not determine assignment values.
For arbitrary desired values $d_i$ and positive $a_i$, the global assignment
utility
\[
 W(i,s)=a_i u_i(s)+d_i-a_i\mathbb E_{L_i}u_i
\]
preserves the given preferences within every role and has conditional value
$d_i$ in role $i$.
\end{proposition}
\begin{proof}
Within each role the displayed function is a positive-affine transformation
of $u_i$. Taking its expectation under $L_i$ gives $d_i$.
\end{proof}
This constructs different global preferences consistent with the conditional
data. It does not permit independently rescaling an already specified global
preference. In particular, exact equality of the game's numerical payoffs and
cardinal equivalence of conditional preferences answer different questions.

For a positive implication fix common monetary receipts, a reference role 0,
and exact automorphisms $t_i$ sending 0 to $i$. Suppose the participant chooses
a reference action after learning the role, and opponent laws are transported
by these same disclosed maps. Exactness gives
$u_i(t_i s)=u_0(s)$ pointwise, hence equal action values and equal optimized
role values for every common reference opponent law. No equation
$t_i t_j=t_k$ is used. One-shot assignment indifference therefore needs no
coherent section. Conversely equal optimized values alone need not identify
payoff rows: menus $\{1,0\}$ and $\{1,-1\}$ have the same optimum.

\begin{proposition}[Invariant recommendations]\label{prop:invariant-recommendations}
Let $G$ be a finite group of exact typed automorphisms of a finite game.
A $G$-invariant profile law gives equal expected receipts to players in the
same player orbit. Averaging any correlated equilibrium over $G$ gives an
invariant correlated equilibrium. The same conclusion applies to a lottery
on Nash profiles through its induced joint action law.
\end{proposition}
\begin{proof}
If $g$ sends $i$ to $j$, invariance and exactness give
$\mathbb E u_j(s)=\mathbb E u_j(gs)=\mathbb E u_i(s)$.
Pushforward by $g$ maps each recommendation/deviation inequality to the
corresponding inequality for the transported actions and recipient. Thus it
preserves correlated equilibrium; convexity permits averaging. A Nash product
law satisfies every correlated-equilibrium inequality, so a mixture of such
laws does as well.
\end{proof}
These are standard change-of-variables and obedience arguments in the
correlated-equilibrium framework of \citet{aumann1974}; see also
\citet{hart1989}. They are not a new equilibrium existence theorem. Averaging does not generally preserve product
laws or the Nash property of the product of marginal averages. In binary
anti-coordination, the equal lottery on the two pure equilibria gives each
player receipt 1, whereas the product of its uniform marginals gives 1/2.
Thus structural symmetry, invariant procedures and equal realized outcomes
remain distinct. No unresolved general fairness characterization is needed
for any theorem below.
